\section{Canonical alphabet and block gaps}\label{sec:alphabet}
Let $E$ be a nontrivial orthogonal representation of the effective group $H$. On $L^2(\mathbb S(E))$ write $V_a f(u)=f(A_a^{-1}u)$ and let $\Pi$ be Haar averaging. All operators below are restricted to $\mathcal H=(I-\Pi)L^2(\mathbb S(E))$, which each $V_a$ preserves. The complexification is used for adjoints and spectra. Define
\begin{equation}\label{eq:optimized-gap}
 r_1(\mathcal A)=\min_{p\in\Delta(\mathcal A)}\norm{P(p)},\qquad
 P(p)=\sum_a p(a)V_a,\qquad g_1^*=1-r_1^2.
\end{equation}
Equivalent orthogonal actions conjugate these operators. Adding a redundant compact presenting factor does not change them. Repeating a letter does not change the set of admissible averages, although it changes the particular uniform law.

\begin{proposition}[Optimization, full support and simultaneous laws]\label{prop:law}
The minimum in \eqref{eq:optimized-gap} is attained. The function $g(p)=1-\norm{P(p)}^2$ is concave and $2$-Lipschitz in $\ell_1$. Its supremum on full-support laws equals $g_1^*$. Every full-support law has positive gap if and only if $g_1^*>0$. If $m=|\mathcal A|$ and $u$ is uniform, then
\begin{equation}\label{eq:uniform-comparison}
 g(u)\ge g_1^*/m.
\end{equation}
More generally, $q(a)\ge\tau p(a)$ for all letters implies $g(q)\ge\tau g(p)$. For finitely many represented types, individual positive optimized gaps can be certified by one common full-support law.
\end{proposition}
\begin{proof}
Unitary norms give $\norm{P(p)-P(q)}\le\norm{p-q}_1$ and $\norm{P(p)}\le1$. Continuity and compactness give attainment, and the square of a norm is convex. Squaring the preceding inequality gives the stated Lipschitz bound for gaps. Approximation of a maximizing law by its mixtures with the uniform law proves the supremum assertion, not attainment by a full-support optimizer.

If $q=\tau p+(1-\tau)p'$ and $f\in\mathcal H$, convexity of the squared norm yields
\[
 \norm{P(q)f}^2\le\tau\norm{P(p)f}^2+(1-\tau)\norm{P(p')f}^2
 \le(1-\tau g(p))\norm f^2.
\]
This proves domination. Uniform $u$ dominates $p/m$ for every $p$. Any full-support law dominates a positive multiple of a maximizing law. Finally a mixture $\sum_i w_ip_i$ dominates $w_ip_i$ for each type's chosen law; a further positive uniform mixture gives full support. Alternatively \eqref{eq:uniform-comparison} applies to all types at once. Optimizers need not be unique.
\end{proof}

Laziness can be priced precisely: if identity is already an allowed letter, adding it with weight $\eta<1$ loses at most the factor $1-\eta$ in this certificate. If it is not a letter, silently inserting idle steps changes the word interface. Nor does zero one-step gap imply zero gap after such a modification.

\subsection{A dual problem and the limit of finite truncation}
For unit $f\in\mathcal H$, let $Q(f)$ be the real Gram matrix
\[
 Q(f)_{ab}=\operatorname{Re}\langle V_af,V_bf\rangle,
 \qquad \mathscr C=\overline{\conv}\{Q(f):\norm f=1\}.
\]
It is a compact convex subset of real positive semidefinite $m$-by-$m$ matrices with diagonal one.
\begin{proposition}[Finite-dimensional Gram dual]\label{prop:gap-dual}
One has
\begin{equation}\label{eq:gap-dual}
 r_1^2=\min_{p\in\Delta_m}\max_{Q\in\mathscr C}p^TQp
      =\max_{Q\in\mathscr C}\min_{p\in\Delta_m}p^TQp.
\end{equation}
If $r_{1,L}$ is the same norm minimum restricted to spherical harmonics of degrees at most $L$, then $r_{1,L}\uparrow r_1$. These finite problems are convex semidefinite programs when their blocks are supplied. They give upper bounds on $g_1^*$, not finite certificates of positive full gap without a tail estimate.
\end{proposition}
\begin{proof}
For fixed $p$, the squared norm is $\sup_f p^TQ(f)p$. Taking the closed convex hull does not change the supremum. The expression is convex in $p$ and affine in $Q$; the compact convex minimax theorem proves \eqref{eq:gap-dual}. This is the usual convex spectral-design mechanism, not a new minimax principle; compare the finite-chain optimization in \cite{BDX}.

The harmonic spaces are invariant under orthogonal substitutions and Haar averaging. Their finite orthogonal sums increase densely. Consequently the truncated norms increase pointwise to the full norm. Take a subsequence of truncated minimizing laws converging to $p_*$. For every fixed $L_0$, continuity on its finite blocks gives $\norm{P(p_*)}_{\le L_0}\le\lim_L r_{1,L}$. Taking the supremum over $L_0$ gives $r_1\le\lim_L r_{1,L}$; the reverse inequality is immediate. The block norm constraints are
\[
 \begin{pmatrix}tI&P_\ell(p)\\P_\ell(p)^*&tI\end{pmatrix}\succeq0
 \quad(\ell\le L),
\]
with invariant subspaces removed. They describe a finite SDP. Convergence alone supplies no effective error bound for the omitted degrees. The set $\mathscr C$ is finite-dimensional but is specified by all unit vectors in an infinite-dimensional space; its dual is not an effective positivity algorithm.
\end{proof}

\subsection{The block invariant}
For $L\ge1$, give every word $w\in\mathcal A^L$ its actual operator $V_w$ and put
\begin{equation}\label{eq:block-gap}
 r_L=\min_{\pi\in\Delta(\mathcal A^L)}\norm{\sum_w\pi(w)V_w},
 \qquad g_L^*=1-r_L^2.
\end{equation}
A block is $L$ ordinary inputs. This definition supplies a test law, not a larger free input symbol or an extra runtime register.

\begin{theorem}[Finite-block expansion]\label{thm:block-gap}
Let $M=m^{-1}\sum_aV_a$ on $\mathcal H$. The following are equivalent:
\begin{enumerate}
\item $g_L^*>0$ for some finite $L$;
\item $\norm{M^L}<1$ for some finite $L$;
\item the spectral radius of $M$ is less than one.
\end{enumerate}
The equivalence remains true with $M$ replaced by the average under any fixed full-support law. Also $r_{L+L'}\le r_Lr_{L'}$. For finitely many types with positive finite-block gaps, a common length and a common full-support word law work for all of them.
\end{theorem}
\begin{proof}
The uniform product law assigns $m^{-L}$ to every word and hence dominates $m^{-L}\pi$ for any law $\pi$. Proposition~\ref{prop:law}, applied to word operators, gives
\begin{equation}\label{eq:block-domination}
 1-\norm{M^L}^2\ge m^{-L}g_L^*.
\end{equation}
A fixed full-support law gives the same argument with the smallest product mass in place of $m^{-L}$. Thus (1) and (2) are equivalent. The spectral-radius formula $\operatorname{spr}(M)=\lim_n\norm{M^n}^{1/n}$ gives (2)$\Leftrightarrow$(3). Concatenating independently chosen word laws proves submultiplicativity. A common multiple of the finitely many good lengths admits a positive gap for every type by repeating each type's good law; a full-support common word law dominates each one. All invariant complements refer to the original effective action.
\end{proof}

This equivalence identifies a genuine difference between norm and spectral radius of a nonnormal action average. It is not a claim that either can be certified by finitely many unaccompanied harmonic calculations. The next sections make the finite numerical calibration explicit, then apply this distinction to a rational alphabet.
